\chapter{Literature Review}

\section{Sea level change}
\sectioncount{Sea level change}

The IPCC's Sixth Assessment Report (AR6) identified sea level rise as ``an existential threat for island nations, low-lying coastal zones, and the communities, infrastructure, and cultural heritage therein'' \citep{cooleyOceansCoastalEcosystems2022}. A rise of 1~m would increase the global population living below mean sea level from between 34--49 million people today to 77--132 million (\citealt{seegerSeaLevelMuch2026}, with uncertainty due to the choice of digital elevation model and population dataset). This figure also belies the variability in local sea level over multiple timescales \citep{woodworthForcingFactorsAffecting2019} --- storm surges and other short-timescale variability are a hazard for communities living in low-elevation zones. The impacts of sea level rise are not only felt in a binary way (i.e.\ land lies either above or below sea level), but through the increased frequency and magnitude of extreme sea level events \citep{boumisCoevolutionExtremeSea2023}. Exposure to these events is highly regional, with vulnerability concentrated in the Global South \citep{hooijerGlobalLiDARLand2021,rentschlerGlobalEvidenceRapid2023,monioudiImpactsSeaLevel2025}. Consequently, the communities most exposed to sea level rise are often also the least economically and institutionally equipped to adapt.~\note{I could also say something here about local sea level change, vertical land movement, gravitational effects, etc.}

Since the 20\textsuperscript{th} century, global mean sea level rise has accelerated from 1.35 ± 0.57 mm yr\textsuperscript{-1} between 1901--1990 to 3.69 ± 0.48 mm yr\textsuperscript{-1} between 2006--2018 \citep{fox-kemperOceanCryosphereSea2021}. The IPCC have high confidence that this acceleration is the result of anthropogenic warming and that future sea level rise is sensitive to future greenhouse gas emissions. AR6 projects a likely (66\% chance) sea level rise of 0.33--0.62 m by 2100 under a low emissions pathway and 0.63--1.01 m under a high emissions pathway, the upper bound of which coincides with the 1 m sea level rise hypothetical evaluated by~\cite{seegerSeaLevelMuch2026}. Sea level rise occurs either through thermosteric effects (``thermal expansion'') or through barystatic contributions (i.e.~mass flux) from another source such as land water storage, mountain glaciers, or the polar ice sheets. The dominant drivers of historic sea level rise have been thermal expansion and mountain glaciers. However, contributions from the ice sheets have accelerated since the 1970s and are projected to be significant sources of sea level rise this century \citep{fox-kemperOceanCryosphereSea2021}. Antarctica is the largest source of uncertainty in these projections, due in large part to poorly constrained processes governing the ice sheet response \citep{edwardsProjectedLandIce2021}.~\note{I'm uncertain about the etiquette for citing the same report multiple times in one paragraph. Have I done this appropriately?}

Beyond 2100, AR6 cites ``deep uncertainty'' in multi-century sea level projections. The 21\textsuperscript{st} century projections are informed heavily by the Ice Sheet Model Intercomparison Project \citep{nowickiIceSheetModel2016}. However, no such project existed for 2300 projections at the time. Instead, AR6 relied on a combination of methods, including extrapolation of 2100 rates of mass loss and the expert elicitation by~\cite{bamberIceSheetContributions2019}. They also considered Marine Ice Cliff Instability (MICI), a highly uncertain process related to the failure of tall ice cliffs, which was only included in a handful of modelling studies at the time, but which dramatically increase Antarctica's mass loss on these timescales \citep{decontoContributionAntarcticaFuture2016,decontoParisClimateAgreement2021}. Even excluding this process, AR6 projected 0.3--3.1~m of sea level rise by 2300 under a low emissions pathway and 1.7--6.8~m under high emissions. Synthesized estimates for the individual sources were not provided, but uncertainty in Antarctica's contribution was wider than 3 m in some assessments of the high emission pathway. Understanding and constraining this uncertainty is therefore a significant goal towards more confident multi-century sea level projections. Since AR6, an increased number of studies have simulated the Antarctic response to multi-century climate change (see \Cref{section:antarctic-multicentury-projections}), but constraining the uncertainty in its sea level contribution remains a challenge, thus demonstrating the need for further work.

% \begin{table}[h]
%     \renewcommand{\arraystretch}{1.5}
%     \centering
%     \caption{Observed land ice sea level contribution and one example of projected contribution, adapted from~\cite{fox-kemperOceanCryosphereSea2021} Tables 9.5, 9.8, and 9.11. Given are the most likely values for the historical period and the medium-confidence likely range (17th-83rd percentile) by 2100 and 2300 under a high emission scenario (relative to a 1995--2014 baseline).}
%     \begin{tabular}{cccc}
%         \toprule
%         Source & 1901--2018 (mm) & 2100 (m) & 2300 (m)  \\
%         \midrule
%         Glaciers    & 67.2  & 0.15--0.20 & 0.32         \\
%         Greenland   & 40.4  & 0.09--0.18 & 0.32--1.75   \\
%         Antarctica  & 6.7   & 0.03--0.34 & -0.28--3.13  \\
%         \bottomrule
%     \end{tabular}\label{table:land-ice-slc}
% \end{table}

\section{The Antarctic Ice Sheet}

\subsection{Geographical setting}
\sectioncount{Geographical setting}

The Antarctic Ice Sheet is the largest store of freshwater on Earth, holding 26.0~±~0.4~million~km\textsuperscript{3} of ice \citep{morlighemDeepGlacialTroughs2020}. It is a dynamic ice mass, accumulating ice through precipitation and flowing outwards under gravity. This flow organises into ice streams, which are fast-flowing corridors of ice that transport ice from the ice sheet interior to its perimeter \citep{bennettIceStreamsArteries2003}. 93\% of the ice sheet perimeter interfaces with the ocean, either as a floating ice shelf (74\%) or vertical ice cliff (19\%) \citep{bindschadlerGettingAntarcticaNew2011}. Ice cliffs discharge ice mainly through the calving of icebergs \citep{bennCalvingGlaciersIce2018} and ice shelves through a combination of calving and basal melting \citep{davisonAnnualMassBudget2023}. The transition point between grounded ice sheet and floating ice shelf is called the ``grounding line'', although in practice this is challenging to define and measure \citep{friedlRemoteSensingGlacier2020}. Ice shelves don't contribute directly to sea level rise, but have indirect consequences for sea level through the buttressing force they impose on upstream grounded ice (see \Cref{section:calving-and-hydrofracture}).

Antarctica can be split into three climatically distinct regions: West Antarctica, East Antarctica, and the Antarctic Peninsula. The West Antarctic Ice Sheet (WAIS) is grounded on bedrock that lies predominantly below sea level (a marine ice sheet), which makes it susceptible to instabilities that could bring about self-sustaining mass loss (see \Cref{section:ice-sheet-instability}). This is particularly concerning because West Antarctica is the region with the highest current rate of mass loss (\Cref{table:component-ice-sheets}) and may have collapsed completely in past warm climates (see \Cref{section:antarctica-in-past-climates}). West Antarctic ice streams drain into the Amundsen Sea, Ross Sea, and Weddell Sea sectors of the Southern Ocean. The latter two feature Antarctica's two largest ice shelves, the Ross ice shelf (\todo{value} km$^2$) and the Filchner-Ronne ice shelf (\todo{value} km$^2$), respectively. The ocean cavities beneath them are key regions for the production of dense deep waters that form part of the global ocean circulation \citep{rintoulAntarcticBottomWater2025}.

\begin{table}[h]
    \renewcommand{\arraystretch}{1.5}
    \centering
    \caption{Antarctica's component ice sheets, their rate of mass loss between 1992--2020 \citep{otosakaMassBalanceGreenland2023a}, and their total sea level potential \citep{morlighemDeepGlacialTroughs2020}.}
    \begin{tabular}{ccc}
        \toprule
        Region & Mass loss (Gt yr$^{-1}$) & Sea level potential (m) \\
        \midrule
        WAIS & -82 ± 9 & 5.3 ± 0.2 \\
        EAIS & 3 ± 15 & 52.2 ± 0.7 \\
        APIS & -13 ± 5 & 0.27 ± 0.02 \\
        \bottomrule
    \end{tabular}\label{table:component-ice-sheets}
\end{table}

The East Antarctic Ice Sheet (EAIS) is separated from West Antarctica by the Trans-Antarctic mountains. It represents by far the largest share of Antarctica's sea level potential, but has contributed the least to recent sea level change (\Cref{table:component-ice-sheets}). Despite the current trends, the EAIS may have been more dynamic in past climates and its future response is highly uncertain (\citealt{stokesResponseEastAntarctic2022}). Although the EAIS is mostly terrestrial (grounded above sea level), its underlying bedrock has three key marine basins: The Recovery, Wilkes, and Aurora subglacial basins. The Totten glacier, which drains most of the Aurora subglacial basin, alone holds 3.89 m sea level equivalent of ice \citep{morlighemDeepGlacialTroughs2020}. These basins may be susceptible to the same instabilities as West Antarctica if they are sufficiently pushed from their current state \citep{mengelIcePlugPrevents2014,aitkenRepeatedLargescaleRetreat2016,golledgeEastAntarcticIce2017}.

The Antarctic Peninsula Ice Sheet (APIS) is a comparatively small part of Antarctica, but is losing mass faster than any other region relative to its size (\Cref{table:component-ice-sheets}). It is distinct as the most Northerly region of the continent, reaching 63$^\circ$S towards the Drake Passage. Its mountain range features smaller, Alpine-like glaciers down steep terrain. The Peninsula acts as an orographic barrier between the Weddell Sea and the Bellinghausen Sea, strongly influencing the climate on either side \citep{daviesAntarcticPeninsulaPresent2026}. The largest ice shelf in the APIS, the Larsen ice shelf, extends down the Eastern flank of the peninsula and is divided into smaller sections Larsen A to D, the largest of which is Larsen C at \todo{value and cite} km$^2$. The recent changes in the Antarctic Peninsula are significant not only to the local environment, but also serve as a bellwether for processes that may later occur in East and West Antarctica with substantially greater consequences for sea level.

\subsection{Recent observed changes}\label{section:recent-observed-changes}

Antarctica is not in equilibrium with the current climate, with some areas showing a clear trend in mass balance that is consistent across multiple methods of measurement (e.g.~\citealt{gardnerIncreasedWestAntarctic2018,smithPervasiveIceSheet2020,velicognaContinuityIceSheet2020}). Furthermore, mass loss from West Antarctica has accelerated across every four year period from -37 ± 19 Gt yr\textsuperscript{-1} in 1992--1996 to -131 ± 21 Gt yr\textsuperscript{-1} in 2012--2016, before slowing to -94 ± 25 Gt yr\textsuperscript{-1} in 2017--2020 \citep{otosakaMassBalanceGreenland2023a}. The vast majority of this mass loss has originated from the Amundsen Sea Sector of West Antarctica, particularly the Thwaites and Pine Island Glaciers \citep{rignotFourDecadesAntarctic2019}. Sediment records from the Amundsen Sea indicate that Thwaites and Pine Island have been in disequilibrium since at least the 1940s \todo{cite}. Direct measurements of Pine Island Glacier made possible by the satellite era reveal a consistent pattern of change: retreat of the glacier terminus, acceleration of its ice stream, and a subsequent thinning as the glacier adjusts to maintain its flux balance \citep{joughinTimingRecentAccelerations2003,payneRecentDramaticThinning2004} \todo{cite more recent}. It is difficult to draw direct links between the current disequilibrium and anthropogenic climate change, but it's likely that climate change has at least exacerbated the retreat of Pine Island glacier over this period \citep{bradleyQuantifyingAttributingRole2025}.~\todo{See Ines's 2023 remote sensing review paper for more inspiration on Pine Island}.

\todo{Thwaites}

Unlike the Arctic, which has warmed by \todo{value} times the global average due to nonlinear feedbacks \citep{previdiArcticAmplificationClimate2021}, Antarctic surface air temperatures have broadly followed the global trend \todo{cite}. The one exception to this is the Antarctic Peninsula, which has warmed three times faster than the rest of the continent, up to temperatures otherwise unprecedented in the Holocene \citep{vaughanRecentRapidRegional2003}. This atmospheric warming has increased surface melting of the Antarctic Peninsula Ice Sheet and its resulting sea level contribution \citep{vaughanRecentTrendsMelting2006}. The most pronounced changes in the APIS have been the thinning, retreat, and collapse of some of the peninsula's ice shelves~\citep{cookOverviewArealChanges2010,rignotIceShelfMeltingAntarctica2013}. The extent of ice shelves on the peninsula roughly trace the -9$^\circ$C isotherm, suggesting a ``limit of viability'' for APIS ice shelves that will shrink with further warming \citep{morrisSpatialTemporalVariation2013}. Although this is a useful heuristic, the mechanisms of ice shelf collapse are not purely thermal (see, for example, Sections~\ref{section:ice-ocean-interactions} and~\ref{section:calving-and-hydrofracture}).

The first ice shelf collapse to occur during the satellite era was the Wordie Ice Shelf in 1989 \citep{doakeRapidDisintegrationWordie1991}. This was followed by a sequential collapse from North to South down the Eastern flank of the peninsula, starting with the Larsen inlet --- also in 1989 --- then the Prince Gustav and Larsen A ice shelves in 1995 \citep{cookOverviewArealChanges2010}. Finally, the Larsen B ice shelf progressively thinned between 1992 and 2001 before its eventual collapse over a few days in February 2002 \citep{shepherdLarsenIceShelf2003}. The Larsen C ice shelf has so far avoided collapse, but has continued to thin at a rate of -3.8 m per decade over increasingly large areas of the ice shelf \citep{paoloVolumeLossAntarctic2015}. Between August 2015 and July 2017, a 200 km crack developed along the ice shelf, leading to the calving of a large iceberg representing 10\% of the total ice shelf area \citep{hoggImpactsLarsenCIce2017}. Due to the loss of buttressing force provided by ice shelves (see Section~\ref{section:ice-ocean-interactions}), the collapse of each ice shelf on the peninsula was consistently followed by an acceleration and thinning of its upstream glaciers \citep{rignotAcceleratedIceDischarge2004,glasserIceshelfTributaryTidewater2011,domgaardHalfCenturyDynamic2025}. The collapse of the Prince Gustav and Larsen A ice shelves cannot be uniquely attributed to anthropogenic climate change, since sediment records indicate that they have previously collapsed during the mid-Holocene \citep{pudseyFirstSurveyAntarctic2001,brachfeldHoloceneHistoryLarsenA2003}. This is corroborated by written accounts of the Prince Gustav ice shelf by 19th century polar explorers suggesting that it had been in retreat since at least 1843 \citep{cooperHistoricalObservationsPrince1997}. However, similar records from ocean sediments previously beneath the Larsen B ice shelf indicate that it had been stable throughout the Holocene, and therefore suggest that its recent collapse contains an anthropogenic signal \citep{domackStabilityLarsenIce2005}.

\note{Too much on the Peninsula? Two big paragraphs on the smallest region (and least important for my work) seems like a lot. On the other hand, the peninsula is an excellent case study for what might happen elsewhere in Antarctica, an argument which I plan to develop in later sections, so it is perhaps worth going into detail here.}

In contrast to the WAIS and APIS, sea level contribution from the East Antarctic Ice Sheet remains close to a state of balance, with a slight positive mass trend of 3 ± 15 Gt yr$^{-1}$ between 1992--2020 \citep{otosakaMassBalanceGreenland2023a}. The high uncertainty in this estimate is the result of disagreement between different methods of measurement, with some studies showing a negative mass trend over similar periods (e.g.~\citealt{rignotFourDecadesAntarctic2019,shepherdTrendsAntarcticIce2019}). These aggregate measurements also belie the spatial variability in East Antarctic mass trends. Several studies highlight regional mass loss from Wilkes Land in East Antarctica, particularly the Totten Glacier, which drains the Aurora subglacial basin \citep{shepherdTrendsAntarcticIce2019,velicognaContinuityIceSheet2020,smithPervasiveIceSheet2020}. This is paired with a positive mass trend over the East Antarctic plateau, consistent with reanalysis-derived increases in precipitation \citep{davisSnowfallDrivenGrowthEast2005}. This spatial variability suggest that whilst the East Antarctic Ice Sheet is in a state of balance, it is not in a state of equilibrium, with the potential for more significant mass change if this balance is disrupted.

\todo{Check Jamie's and Mira's theses to look for more EAIS notes.}

\todo{Quick synthesis of this section?}

\subsection{Antarctica in past climates}\label{section:antarctica-in-past-climates}
\sectioncount{Antarctica in past climates}

\note{Comparing this to Jamie's section on Antarctica in the Pliocene, mine is much more condensed and missing quite a lot of context. However, I wouldn't expect my section to be as big as Jamie's anyway, given the relative importance of the Pliocene in my thesis. For example, I essentially pick up with Pliocene sea level estimates in 2015, whereas Jamie gives a longer history. I should probably include some discussion of palaeoshoreline work, which gave even higher sea level estimates than d18O. I could also spend time explaining how the $\delta^{18}$O record works. I generally feel that I am rushing through references without giving them proper explanation or room to breath, but this may also be a matter of taste.}

Past climates are a natural laboratory for exploring the how the Earth system responds to changes in greenhouse gases. By comparison with the modern era, they can inform projections of the future climate \citep{tierneyClimatesInformOur2020}. The Mid-Pliocene Warm Period (mPWP), a period 3.3--3.0 million years ago with CO\textsubscript{2} concentrations near 400 ppm and temperatures 3--5.3$^\circ$C warmer than present \citep{tierneyPlioceneWarmthPatterns2025}, is the closest analogue to end of the century projections under medium to high emissions \citep{burkePlioceneEoceneProvide2018}.~\note{I've used the Tierney 2025 numbers here because they seem quite convincing, but perhaps I should stick to the IPCC values and simply mention that it could be warmer?} Early attempts to estimate Pliocene sea level based on the the $\delta^{18}$O record --- preserved in the layers of benthic foraminifera that occupy ocean sediment cores --- produced highly uncertain results \citep{duttonSealevelRiseDue2015}. For example,~\cite{millerHighTideWarm2012} estimated that sea level was 22~±~10~m higher than today, the lower bound of which would require the near total loss of the Greenland and West Antarctic Ice Sheets \note{This is just technically false because miller et al use a multiproxy approach, not just d18O}. The upper range of these estimates then requires significant contribution from marine and possibly terrestrial sectors of East Antarctica. Much lower sea level estimates can be obtained by using a higher $\delta^{18}$O of Antarctic ice, justifiable by warmer temperatures of the Pliocene \citep{winnickOxygenIsotopeMassbalance2015,gassonModelingOxygenIsotope2016}, but the true value is then highly uncertain.~\cite{raymoAccuracyMidPlioceneD18Obased2018} highlight further biases in the $\delta^{18}$O record that could both under-~and over-estimate Pliocene sea level.

\note{I've again condensed a lot of information below compared to Jamie's. The references I've cited could all be explained in much more detail, but perhaps this level of detail is appropriate for my thesis?}

Regional palaeoclimate records provide a suite of evidence for a dynamic WAIS that advanced and retreated in obliquity-paced glacial-interglacial cycles:~\cite{naishObliquitypacedPlioceneWest2009} and~\cite{seidensteinPliocenePleistoceneWarmwater2024} found open-ocean diatom species and warm-water foraminifera in sediments beneath the current Ross ice shelf;~\cite{gohlSeismicStratigraphicRecord2013,gohlEvidenceHighlyDynamic2021} reveal buried grounding zone wedges in various locations of the stratigraphic record; and~\cite{passchierWestAntarcticIce2025} and~\cite{horikawaRepeatedMajorInland2026} demonstrate a cyclicity in the sediment composition of the continental shelf. In contrast, evidence for change in the EAIS is more fragmentary, but an increasing number of palaeoclimate records point to periods of substantial retreat. Sediment cores offshore of the George V and Adélie coasts show a sequence of change in ocean productivity, ice-rafted debris, and sediment provenance pointing to retreat in the Wilkes subglacial basin \citep{cookDynamicBehaviourEast2013,pattersonOrbitalForcingEast2014,bertramPlioceneDeglacialEvent2018}. Ice rafted debris from Wilkes land and Adélie land has been found as far away as Prydz Bay, which would require large icebergs to transport, perhaps comparable to Northern Hemisphere Heinrich events \citep{williamsEvidenceIcebergArmadas2010}. This is supported by deep glacial erosion representing a plausible location for the terminus of Totten Glacier in the Pliocene, 150 km behind its current grounding line \citep{aitkenRepeatedLargescaleRetreat2016}. Conversely, seismic soundings of the Aurora subglacial basin evidence the existence of ice near current grounding lines since the late Miocene \citep{gulickInitiationLongtermInstability2017}. The EAIS also appears to have responded differently to orbital forcing than the WAIS, lacking a clear obliquity signal in its climate records \citep{pattersonSpatiallyVariableResponse2026}.

\note{Again, the below paragraph could go into much more detail, particularly about PLISMIP.\@ I'm interested to know if you feel any particular paragraph in this section is too short, or whether they are all too short (or maybe not.)}

The regional records provide coherent evidence of significant mass loss from the WAIS and at least some of the marine sectors of the EAIS.\@ However, ice sheet models in the Pliocene Ice Sheet Model Intercomparison Project \citep{dolanPlioceneIceSheet2012} were unable to reproduce this mass loss when forced by a Pliocene climate \citep{deboerSimulatingAntarcticIce2015,dolanHighClimateModel2018}. Most (but not all) models simulated a collapse of the WAIS, but this was often offset or even exceeded by accumulation over East Antarctica. The theory of marine ice cliff instability (\Cref{section:ice-sheet-instability}) was first popularised by its ability to dramatically increase Antarctic sea level contribution in the mPWP \citep{pollardPotentialAntarcticIce2015}. However, it was not required to reach the lower range of Pliocene sea level when accounting for modelling and reconstruction error \citep{edwardsRevisitingAntarcticIce2019}. Other models were later able to simulate higher mass loss from Antarctica without MICI, including from the marine sectors of East Antarctica \citep{golledgeAntarcticClimateIcesheet2017,blascoAntarcticTippingPoints2024}. Furthermore, a grain size proxy that avoids the high uncertainties related to $\delta^{18}$O suggests that sea level oscillated by 4.1--20.7 m over Pliocene glacial-interglacial cycles \citep{grantAmplitudeOriginSealevel2019,grantPlioceneSeaLevel2021}. If taken as a deviation from present-day sea level, the upper bound would still require mass loss from marine sectors of East Antarctica, but not from terrestrial sectors.

This progress in both models and palaeoclimate records has made significant progress towards closing the gap between them. Both support a dynamic WAIS that collapsed in climates not so different from that projected by the end of this century. Marine sectors of the EAIS may also be vulnerable to anthropogenic warming, but the palaeoclimate records cannot tell us anything about the pace of retreat --- whether it would occur on multi-century timescales. The uncertainties in sea level reconstructions are still very high, but the changes to Antarctica are very large, making them a natural complement to observations in the modern era (with small uncertainties but small changes) when assessing the results of ice sheet models.

\section{Ice sheet dynamics}
\subsection{Ice rheology*}

\note{I may want a section here to talk about the polycrystalline structure of ice, as this is relevant to the flow law and the uncertain value of Glen's exponent, but perhaps this is getting too into the weeds.}

\subsection{Ice-atmosphere interactions*}

\note{Not particularly relevant to my thesis given that we take SMB directly from GCMs. However, a quick review of lapse rates might be useful when discussing why we leave them out later. Anything else in ice-atmosphere that would be good to talk about?}

\subsection{Ice-ocean interactions}\label{section:ice-ocean-interactions}
\sectioncount{Ice-ocean interactions}

The Southern ocean is a complex system of different water masses that occupy certain depths in the water column, varying significantly in both space and time. We refer the reader to~\cite{wangWaterMassesSouthern2025} for a full review, but we highlight several important water masses below that are significant for the Antarctic Ice Sheet.

The uppermost layer of the Southern Ocean, referred to as Antarctic Surface Water (AASW) is governed by turbulent mixing, driven mainly by wind stresses at the surface. The ocean surface is very close to the local freezing point of sea water, approximately -1.8$^\circ$C. These temperatures circulate through the mixed layer, which varies seasonally in thickness due to changes in the winds. The mixed layer is deepest in winter, reaching depths of 100--200 m or more, but retreats in summer, leaving behind a cold layer known as Winter Water (WW). This layer maintains its cold temperatures whilst the surface waters above it warm in summer. The salinity of AASW is between \todo{value} psu, less than the global average of \todo{value} psu. This is for two reasons: firstly, the Southern ocean has a net positive freshwater exchange with the atmosphere, due to receiving more freshwater in precipitation than it loses in evaporation; secondly, it receives freshwater from the melting of sea ice and the Antarctic ice sheet. This low-salinity is inherited by WW, though summer surface water is fresher due to that year's summer ice melt.

The Southern Ocean is also a confluence of deep water masses that are transported through the global overturning circulation \citep{talleyClosureGlobalOverturning2013}. Deep water from the Pacific, Indian, and North Atlantic Oceans mix to form Circumpolar Deep Water (CDW). From these water masses, CDW inherits warm temperatures of \todo{value} $^\circ$C and high salinity of \todo{value} psu. It occupies the water column at depths of up to 3000 m, but shoals further South due to regional upwelling. The prevailing Westerly winds that circle the Antarctic continent create a meridional Ekman flow, driving surface waters equatorward and drawing up deep waters to replace them. In these regions, CDW sits just below WW, forming a steep temperature and salinity gradient between the two water masses that appears in ocean depth profiles.




deepest in the winter due to intense winds, reaching depths of 100--200 m or more. This creates the uppermost layer of the Southern Ocean, referred to as Winter Water (WW).

is occupied by very cold surface waters, with temperatures close to the local freezing point of sea water, approximately -1.8$^\circ$C. This water tends also to be very fresh, due to the freshwater fluxes from sea ice and the Antarctic ice sheet. This maintains a degree of stratification

The Southern Ocean is a confluence of water masses transported through the global overturning circulation. Deep water from the Pacific, Indian, and North Atlantic Oceans mix to form Circumpolar Deep Water (CDW) at depths of up to 3000m \citep{talleyClosureGlobalOverturning2013}. CDW is characterised by warm temperatures and high salinity 


which circulates through the Antarctic Circumpolar Current (ACC) at depths of up to 3000 m \citep{talleyClosureGlobalOverturning2013}. CDW is characterised generally by much warmer temperatures and higher salinity than the surface waters above it \citep{wangWaterMassesSouthern2025}. It can be divided into `upper' and `lower' water masses with distinct characteristics. Upper CDW includes contributions from Indian and Pacific deep waters, which are warm, oxygen-poor, and nutrient-rich; Lower CDW contains more North Atlantic Deep Water, giving it characteristically high salinity \citep{whitworthWaterMassesMixing2013}. Upper CDW shoals as it moves Southwards, eventually mixing with surface waters at the Southern boundary of the ACC.\@ However, lower CDW is dense enough to remain below the mixed layer and extend further towards the Antarctic continent \citep{orsiMeridionalExtentFronts1995}. The melting of Antarctic ice shelves is predominantly a question of whether this CDW is able to access the continental shelf that surrounds Antarctica \citep{wangMultiscaleOceanicHeat2023}.

Access to the shelf is limited by the Antarctic Slope Front (ASF), a density front along the shelf break, and a corresponding westward  Antarctic Slope Current (ASC) \citep{jacobsNatureSignificanceAntarctic1991}. The ASF acts as a barrier between CDW and the cooler, fresher shelf waters that occupy ice shelf cavities around Antarctica. The structure of the ASF is highly variable across different sectors of Antarctica and is sensitive to local wind stresses \citep{spenceRapidSubsurfaceWarming2014} and freshwater fluxes \citep{siHeatTransportAntarctic2023}. These local conditions give rise to three distinct ice shelf cavity regimes: cold and fresh cavities, cold and saline cavities, and warm and saline cavities.

In Warm and saline cavities such as those in the Amundsen Sea, a weak ASF allows CDW to encroach onto the shelf through 

The ASF is weakest in the Amundsen Sea, where the Eastward Antarctic Circumpolar Current dominates over the westward ASC \citep{}. This Eastward flow is associated with a polewards Ekman transport in the layer below, which helps to transport CDW onto the continental shelf \citep{klinckExchangeShelfBreak2010}. This oc

The Southern boundary of the ACC flows very close to the continental shelf in the Amundsen and Bellinghausen Seas, but is deflected away from the Antarctic continent by cyclonic gyres in the Weddell and Ross Seas \citep{orsiMeridionalExtentFronts1995}. Temperatures in the gyres are generally lower than those found in ACC \citep{wangMultiscaleOceanicHeat2023}, but some heat is transported to the shelf break along the Eastern boundary of each gyre via entrained CDW \citep{assmannVariabilityDenseWater2005,ryanWarmInflowEastern2016}. 

have strengthened and shifted polewards as a response to greenhouse gas emissions and ozone depletion \citep{waughRecentChangesVentilation2013,leeDetectingOzoneGreenhouse2013}.

\subsection{Buttressing, calving and hydrofracture*}\label{section:calving-and-hydrofracture}
\note{Placemarker section in case I have time to discuss this in more detail.}

\subsection{Subglacial hydrology*}
\note{Placemarker section in case I have time to discuss this in more detail.}

\subsection{Ice sheet instability}\label{section:ice-sheet-instability}

Since the 1970s, it has been suggested that the West Antarctic ice sheet may undergo rapid deglaciation in a warming climate \citep{mercerWestAntarcticIce1978}. This has evolved into the theory of Marine Ice Sheet Instability (MISI), which states that the grounding line position of a marine ice sheet on a retrograde slope (i.e. the bed becomes deeper towards the centre) is inherently unstable \citep{schoofIceSheetGrounding2007}. The fundamental principle of MISI is that discharge at a glacier's terminus is proportional to the glacier thickness at its grounding line. The retreat of the grounding line down a retrograde slope increases thickness, which in turn increases discharge, which causes the grounding line to retreat further, creating a feedback loop. In practice, this effect is mitigated by ice shelves, which buttress the ice sheet, generating sufficient friction to stabilise the grounding line \citep{joughinStabilityWestAntarctic2011}. However, Antarctica’s ice shelves are also thinning as a result of climate change. In 2002  the Larsen A and B ice shelves on the Antarctic peninsula collapsed over just a few days \citep{shepherdLarsenIceShelf2003}.

The West Antarctic Ice Sheet is a strong candidate for undergoing MISI, since it overlies a deep bathymetric basin up to 2000m below sea level, giving it an overall retrograde slope. This is supported by satellite measurements of Antarctic mass balance, which identify the Amundsen Sea Embayment in West Antarctica as the primary region for ice loss in recent decades \citep{smithPervasiveIceSheet2020,velicognaContinuityIceSheet2020,otosakaMassBalanceGreenland2023a}. Pine Island glacier, which terminates in the Amundsen Sea, shows evidence of Circumpolar Deep Water at its grounding line, retreat down a retrograde slope, thinning of its ice shelves, and an acceleration in its discharge \citep{joughinTimingRecentAccelerations2003,jenkinsObservationsPineIsland2010,jacobsStrongerOceanCirculation2011,favierRetreatPineIsland2014}. It has been proposed that West Antarctica could be approaching a 'tipping point', in which MISI could push the ice sheet into a dramatically new state with only a small change in forcing \citep{pattynUncertainFutureAntarctic2020,rosierTippingPointsEarly2021,armstrongmckayExceeding15degCGlobal2022}. Regional ocean models that are able to resolve the temperature in ice shelf cavities find that further warming is unavoidable under any future scenario \citep{naughtenUnavoidableFutureIncrease2023}, indicating that collapse of the West Antarctic ice sheet may now be inevitable.

Compared to the WAIS, the EAIS has been broadly considered to be relatively stable. However, the EAIS, whilst mostly above sea level, does include some marine sectors which may be vulnerable to MISI \citep{armstrongmckayExceeding15degCGlobal2022}. As with the Amundsen Sea Embayment, warm Circumpolar Deep Water has been observed infiltrating the ice shelf cavities and driving thinning of the Totten glacier, which drains part of the Aurora subglacial basin \citep{greenbaumOceanAccessCavity2015,aitkenRepeatedLargescaleRetreat2016,rintoulOceanHeatDrives2016,silvanoDistributionWaterMasses2017}. The ice shelf at the terminus of the Cook glacier, which drains a large proportion of the Wilkes basin, collapsed between 1973-1989, and basal thawing has recently been observed along the nearby Adélie-George V coast \citep{milesVelocityIncreasesCook2018,dawsonHeterogeneousBasalThermal2024}. However, warming over the Southern ocean also drives increased precipitation over East Antarctica \citep{nicolaRevisitingTemperatureSensitivity2023}, which has the potential to offset some or all of the mass loss in the marine sectors. Uncertainty in the overall mass balance is large, but a synthesis of multiple modelling and observation studies suggests that it will remain approximately in balance this century, but has the potential to cause multiple metres of sea level rise on multi-century timescales \citep{stokesResponseEastAntarctic2022}.

Other than marine ice sheet instability, additional mechanisms have been proposed that could drive rapid deglaciation in marine segments of Antarctica. Marine Ice Cliff Instability (MICI; \citealt{bassisUpperLowerLimits2012}) is the theory that marine-terminating ice cliffs could be liable to rapid retreat under certain conditions. Models that include MICI yield significantly greater sea level rise than those without \citep{decontoContributionAntarcticaFuture2016,decontoParisClimateAgreement2021}. However, the inclusion of visco-elastic effects in such models indicate that ice cliffs could be more stable than previously thought \citep{clercMarineIceCliff2019,bassisTransitionMarineIce2021}. Bathymetric grooves in the ocean floor could hint at evidence of MICI \citep{wiseEvidenceMarineIcecliff2017}, but it has yet to be observed in the modern era, nor is it necessary to achieve the sea level change seen in past climates \citep{edwardsRevisitingAntarcticIce2019}. Similarly, corrugation ridges along the ocean floor suggest a 'buoyancy-driven' mechanism for ice sheet retreat, up to hundreds of metres per day along low-gradient slopes \citep{batchelorRapidBuoyancydrivenIcesheet2023}. However, like MICI, it has yet to have been observed in the modern climate.

\subsection{Solid Earth effects}

On short to medium timescales, ice-ocean interactions and ice sheet stability are the main dynamics to consider regarding future sea level rise. However, on longer timescales spanning multiple centuries or millennia, additional effects relating to the interaction between the ice sheets, the solid Earth, and its gravitational field, become relevant. These effects are collectively referred to as Glacial Isostatic Adjustment (GIA; \citealt{whitehouseGlacialIsostaticAdjustment2018}, and references therein). As an ice sheet retreats, the load that it exerts on the Earth’s crust decreases, inducing a visco-elastic response which uplifts the land surface due to upwards pressure in the asthenosphere. GIA is currently observed in regions that experienced deglaciation following the last ice age, recording uplift of over 10 mm a\textsuperscript{-1} in some parts of North America and Europe \citep{sellaObservationGlacialIsostatic2007,lidbergRecentResultsBased2010}. As evidenced by these regions, which have been without an ice sheet for over 8,000 years, the response to glacial unloading is not instantaneous and can be drawn out for millennia. 

GIA has the capacity to significantly influence shoreline migration and bedrock topography over long timescales, which in turn feeds back into the dynamics driving ice sheet growth or retreat. For example, the uplift of bedrock following unloading could limit the accessibility of Circumpolar Deep Water or reverse/intensify the retrograde slope in a region susceptible to Marine Ice Sheet Instability. \cite{larourSlowdownAntarcticMass2019} find that GIA mitigates 27\% of projected sea level rise over 500 years using a model of the Thwaites glacier in which it is included. Similarly,~\cite{kachuckRapidViscoelasticDeformation2020} show that GIA could reduce ice mass loss from the Pine Island Glacier by 30\% over 150 years. These modelling studies show that GIA can have an impact on even sub-millennial timescales.

\section{Ice sheet modelling}
\sectioncount{Ice sheet modelling}

Numerical ice sheet models are important tools for Antarctic science, both for understanding the mechanisms behind recent ice sheet change (e.g.~\citealt{payneRecentDramaticThinning2004,gudmundssonInstantaneousAntarcticIce2019,bradleyQuantifyingAttributingRole2025}), exploring possible ice sheet configurations in the deep past (e.g.~\citealt{deboerSimulatingAntarcticIce2015}) and for making projections of future ice sheet change (\citealt{seroussiISMIP6AntarcticaMultimodel2020,seroussiEvolutionAntarcticIce2024}, among many others). This section reveiws the construction of a general ice sheet model and the important choices involved in simulating the Antarctic ice sheet.

\subsection{Stokes flow}
\sectioncount{Stokes flow}

\noindent
Practical ice flow modelling is done using a hierarchy of approximations to incompressible Stokes flow, which is itself the low Reynolds number limit of the Navier-Stokes equations. In this limit, inertial terms are ignored, and flow is governed by a balance between pressure gradients, viscous stresses, and gravitational driving stresses. These are solved alongside a constitutive relation, which defines the ice rheology (see Section~\ref{sec:rheology}). It is also necessary to impose boundary conditions at the base and surface of the ice sheet. For a rigid boundary (the interface at the bed), this might take the form:

\begin{equation}
    \mathbf{u}_b = F (\boldsymbol{\tau}_b), \quad \mathbf{u}.\mathbf{n} = -m_b,
\end{equation}

\noindent
where $\mathbf{u}$ is the ice velocity field, $\mathbf{u}_b$ is the velocity component parallel to the bed, $\mathbf{n}$ is the normal to the bed, $\boldsymbol{\tau}_b$ is the basal shear stress, $F$ is a sliding law (see Section~\ref{sec:sliding}), and $m_b$ is the basal melt rate (in velocity units). For a free surface (the interface with the atmosphere or ocean), dynamic and kinematic boundary conditions are applied to specify the boundary stresses and surface evolution, respectively. Some model setups include an advection-diffusion equation for temperature, which couples to the momentum (Stokes) equations through the constitutive relation (Eq.~\ref{eq:glenslaw}). However, assuming isothermal ice (i.e.~fixed temperature) is often a reasonable approximation over short timescales.

\subsection{Constitutive relation}\label{sec:rheology}
\sectioncount{Constitutive relation}

The constitutive relation describes the relationship between the stress and deformation of a material. For ice, it is common to use Glen's flow law \citep{glenCreepPolycrystallineIce1955}:
\begin{equation}
    \label{eq:glenslaw}
    \dot{\epsilon}_{ij} = A \tau^{n-1} \tau_{ij}, \quad \tau = \sqrt{\frac{1}{2}\tau_{ij}\tau_{ij}},
\end{equation}

\noindent
where $\epsilon_{ij}$ is the strain rate tensor, $\tau_{ij}$ is the deviatoric stress tensor, $A(T)$ is a temperature-dependant rate factor, and $n$ is a scalar exponent (``Glen's exponent''). Early work to estimate the value of $n$ for polycrystalline ice suggested that $n \approx 3$ (\citealt{weertmanCREEPDEFORMATIONICE1983} and references therein). Subsequently, $n=3$ became standard across a wide range of ice sheet models (e.g.~\citealt{larourContinentalScaleHigh2012,cornfordAdaptiveMeshFinite2013,lipscombDescriptionEvaluationCommunity2019}). However,~\cite{goldsbySuperplasticDeformationIce2001} demonstrated that creep can be divided into multiple regimes, occurring through different mechanisms within the polycrystalline structure. The appropriate value of $n$ depends on the dominant creep mechanism, which in turn depends on the stress.~\note{I could talk here about typical driving stresses in Antarctica and the corresponding value $n=1.8$, but my attempts at this have gotten a bit bogged down and wordy, so have removed for better flow}. Other studies suggest that observations of Antarctic ice streams and ice shelves may be better explained by a higher value of $n$ (e.g.~\citealt{cuffeyHowNonlinearCreep2011,millsteinIceViscosityMore2022}). Ice sheet models that have explored this found Antarctic sea level contribution to be sensitive to the choice of $n$, with higher values leading to greater sea level rise \citep{rosierCalibratedSeaLevel2025,getraerIncreasingGlenNye2025,martinRatesSeaLevelRise2026}.~\note{Is signposting like ``We will explore this uncertainty in chapter 2'' (or similar) appropriate in these situations?}

\todo{Cuffey and Paterson and have a good section on this --- worth a brief rewriting after reading}

\subsection{Approximations in ice flow models}
\sectioncount{Approximations in ice flow models}

Full Stokes modelling is computationally expensive and so is rarely deployed for continent-scale ice sheet modelling (with some notable exceptions, e.g.~\citealt{seddikSimulationsGreenlandIce2012}). Approximations are typically applied to match the requirements of a particular ice-dynamical regime.

\subsection{Basal sliding}\label{sec:sliding}
\subsection{Surface mass balance*}
\subsection{Sub-shelf melting}
\subsection{Glacial isostatic adjustment}

\section{Projections of the Antarctic ice sheet}
\subsection{ISMIP6}
\sectioncount{ISMIP6}

The Coupled Model Intercomparison Project Phase 6 (CMIP6; \citealt{eyringOverviewCoupledModel2016}) is a coordinated set of climate experiments that informed the projections of future climate in AR6. Its subsidiary, the Ice Sheet Model Intercomparison Project for CMIP6 (ISMIP6; \citealt{nowickiIceSheetModel2016}), specifically explored changes in the ice sheets and their contribution to sea level rise. ISMIP6 brought together multiple ice sheet modelling communities to conduct experiments under a strict protocol, such that the results could be directly compared and aggregated across numerous ice sheet models \citep{nowickiExperimentalProtocolSea2020}. This included experiments that explored the ice sheet response to Representative Concentration Pathways, or RCPs \citep{vanvuurenRepresentativeConcentrationPathways2011}. RCPs defined the projected greenhouse gas concentration of the Earth's atmosphere up to 2100 and were an important feature of the previous assessment report \citep{collins2013long}. They were used as boundary conditions for climate models taking part in CMIP5 \citep{taylorOverviewCMIP5Experiment2012}, which simulated the Earth's climate response to these elevated greenhouse gas concentrations. However, climate models at the time invariably included only a static ice sheet in order to conserve computational power. Instead, ISMIP6 used the climate variable outputs of six CMIP5 models as boundary conditions for standalone ice sheet models. The climate models were too coarse-resolution to resolve ice shelf cavities, so ocean properties were extrapolated into the cavity as described in~\cite{jourdainProtocolCalculatingBasal2020b}. The ISMIP6 protocol also included a control simulation, computing sea level contribution under a fixed climate, which was then subtracted from the sea level contribution under each RCP to account for model biases.

The results of ISMIP6 yielded a 2100 Antarctic sea level contribution of between -7.8--30 cm under very high emissions (RCP8.5) and -1.4--15.5 cm under low emissions (RCP2.6) \citep{seroussiISMIP6AntarcticaMultimodel2020}. Higher warming therefore widened the range of sea level contribution in both directions, with a higher maximum and lower minimum. This highlights an important feature of Antarctic 21\textsuperscript{st} century projections, which is the inseparability of emissions scenarios. Many possible futures for Antarctic sea level contribution are reachable from either emission pathway, implying a certain degree of scenario-independence on these timescales. This has been noted by several subsequent studies and cited as important motivation for longer-timescale Antarctic sea level projections \citep{edwardsProjectedLandIce2021,lowryInfluenceEmissionsScenarios2021}. It's also important to note that these results are relative to the control experiment, which itself yielded sea level contribution between -13--14.2 cm, comparable in range to the RCPs. Since the control experiment is under fixed climate, it explores the influence of factors outside of the ice-dynamic response to forcing. This includes the influence of the ice sheet initial state, structural uncertainty in the ice sheet model, uncertainty in the observations, and integration of observational data into the models. These effects were explored further in the InitMIP experiment carried out by ISMIP6 \citep{seroussiInitMIPAntarcticaIceSheet2019}.

The ISMIP6 ensemble also highlights the competing effects of surface mass gains and dynamic mass loss, the balance of which determines the sign of Antarctica's sea level contribution. Under high emissions, ice shelf thinning drives a dynamic response in marine-terminating glaciers, increasing the rate of mass loss. However, climate warming drives an increased surface mass balance, predominantly over East Antarctica, which offsets or exceeds this mass loss. These competing effects drive a significant difference between the WAIS and EAIS response --- the marine WAIS is significantly impacted by ice shelf thinning and therefore generally loses mass, whereas the EAIS experiences both dynamic mass loss in Wilkes Land and surface mass gain over its terrestrial sectors, generally gaining mass overall.~\cite{seroussiInsightsVulnerabilityAntarctic2023} also show that uncertainty in projected Antarctic mass loss is highly regionalised. The response of the Whillans and MacAyeal ice streams, which flow into the Ross ice shelf, is dominated by the choice of climate model forcing, whereas the Thwaites and Pine Island glacier response is most dependant on the choice of ice sheet model and sub-shelf melt parameterisation.

\subsection{Multi-century projections}\label{section:antarctic-multicentury-projections}
\sectioncount{Multi-century projections}

As explained in the introduction to this chapter, ISMIP6 did not make projections of sea level contribution beyond 2100, with multi-century projections relying on extrapolation, expert elicitation, and disparate modelling studies to make the AR6 assessment \citep{fox-kemperOceanCryosphereSea2021}. This is largely due to a lack of appropriate forcing, since state of the art climate model simulations did not generally extend past 2100 at this time. The RCPs were extended to 2500 as described in~\cite{meinshausenRCPGreenhouseGas2011}. However, multi-century projections in CMIP5 were reliant on models of intermediate complexity (\citealt{zickfeldLongTermClimateChange2013}). These were less sophisticated than the more complex General Circulation Models (GCMs) and do not output the necessary climate variables to run an ice sheet model. Multi-century ice sheet projections then either fixed the climate after 2100 \citep{lowryInfluenceEmissionsScenarios2021,chambersMassLossAntarctic2022}, used an idealised forcing \citep{martinMillennialScaleVulnerabilityAntarctic2019,garbeHysteresisAntarcticIce2020}, or constructed an appropriate forcing through statistical methods \citep{golledgeMultimillennialAntarcticCommitment2015,bulthuisUncertaintyQuantificationMulticentennial2019,decontoParisClimateAgreement2021,greveFutureProjectionsAntarctic2023}. For example,~\cite{golledgeMultimillennialAntarcticCommitment2015} and~\cite{bulthuisUncertaintyQuantificationMulticentennial2019} extend the 2100 CMIP5 GCM projections by scaling them linearly with temperature, which is derived from the CMIP5 intermediate complexity models. This provides a proxy for RCP-based multi-century climate simulations, but assumes that the relationship between temperature and other climate variables stays constant with time. Nonetheless, these studies highlight the importance of multi-century forcings, with~\cite{greveFutureProjectionsAntarctic2023} finding roughly twice the sea level contribution of its fixed-forcing equivalent in~\cite{chambersMassLossAntarctic2022} due to continued warming past 2100.

Several advances have been made since ISMIP6 that have improved multi-century projections. They are: new and more policy-relevant emission scenarios, new components added to climate and ice sheet models, and the first multi-century simulations by state of the art climate models. The RCPs were replaced in CMIP6 by the Shared Socioeconomic Pathways (SSPs), which each describe a narrative of how society develops over the 21\textsuperscript{st} century \citep{riahiSharedSocioeconomicPathways2017}. The SSPs correspond roughly with the RCPs in terms of their greenhouse gas concentrations and radiative forcing, but more directly tie into themes of climate policy and international cooperation. These were then extended to 2500 using a decreasing linear trend in fossil fuel emissions, reaching zero by 2250 at the latest \citep{meinshausenSharedSocioeconomicPathway2020}. General Circulation Models have received new components, for example in atmospheric chemistry or the land surface, now referred to as Earth System Models (ESMs). Some models now include interactive ice sheets, allowing them to capture ice-climate feedbacks \citep{muntjewerfDescriptionDemonstrationCoupled2021,siahaanAntarcticContribution21stcentury2022,goelzerInteractiveCouplingGreenland2025}. However, this is still a rarity among Earth System Models, is very expensive to run for multiple centuries, and for this reason is generally not used to explore uncertainty. Most importantly, several climate models (without coupled ice sheets) have now ran climate simulations up to 2300 using the extended SSPs. This has, for the first time, allowed multi-century ice sheet projections without the need for statistical extrapolation of climate forcing.

~\cite{seroussiEvolutionAntarcticIce2024} conducted an ISMIP6-style experiment designed for multi-century projections of the Antarctic Ice Sheet. The experimental protocol included forcing from two climate models with SSP-based projections to 2300: CESM2--WACCM and UKESM1--0--LL.\@ This was accompanied by other models with only 21\textsuperscript{st} century projections, but with forcing extended by repeatedly shuffling the forcing in years from 2080--2100. The experiment also focused mainly on a high emissions future scenario, with only 2 out of 14 recommended forcings representing a low emissions pathway. The results highlight the acceleration of mass loss from Antarctica that occurs after the end of the 21\textsuperscript{st} century, reaching up to 4.4 m sea level equivalent by 2300. However, the lower bound under high emissions was -0.6 m, emphasising the deep uncertainty that remains in multi-century Antarctic sea level contribution. Similar to ISMIP6, the parameterisation of sub-shelf melt and the choice of climate forcing were significant drivers of this uncertainty.



\subsection{Uncertainty Quantification}

\section{Uncertainty quantification*}
\subsection{Experiment design}
\subsection{Gaussian Process Emulation}
\subsection{Calibration}